% =====================================================================
%  CHƯƠNG 6. TAXONOMY ỨNG DỤNG DOANH NGHIỆP, RỦI RO VÀ QUẢN TRỊ AN TOÀN
% =====================================================================
\chapter{Taxonomy ứng dụng doanh nghiệp, rủi ro và quản trị an toàn}\label{ch:quantri}

\section{Phân loại tác tử doanh nghiệp theo Google Cloud và McKinsey}
Google Cloud \cite{google_agents} và McKinsey \cite{mckinsey_agents} nhận diện sáu nhóm tác tử chủ lực trong doanh
nghiệp (Hình~\ref{fig:ch6-taxonomy}):
\begin{itemize}
  \item \textbf{Customer Agents:} chăm sóc khách hàng 24/7, tra cứu đơn hàng, hỗ trợ khiếu nại và đổi trả theo quy chế.
  \item \textbf{Employee Agents:} hỗ trợ nhân sự nội bộ, giải đáp quy trình, đăng ký nghỉ phép, đón nhân viên mới
        (onboarding).
  \item \textbf{Data Agents:} phân tích dữ liệu kinh doanh, chuyển câu hỏi ngôn ngữ tự nhiên thành truy vấn SQL/Cypher và
        xuất báo cáo tự động.
  \item \textbf{Code Agents:} sinh mã nguồn, rà soát lỗ hổng bảo mật, viết kiểm thử đơn vị, hỗ trợ CI/CD.
  \item \textbf{Creative Agents:} thiết kế thông điệp tiếp thị đa kênh, cá nhân hóa nội dung theo hồ sơ khách hàng.
  \item \textbf{Security Agents:} giám sát nhật ký an ninh mạng, phát hiện tấn công bất thường và kích hoạt quy trình ứng
        phó.
\end{itemize}

\hinhtikz{ch6_taxonomy}{Sáu nhóm tác tử doanh nghiệp theo Google Cloud và McKinsey.}{ch6-taxonomy}

\section{Thách thức kỹ thuật và bốn nguyên tắc quản trị theo IBM Think}
Bên cạnh lợi ích, IBM Think \cite{ibm_agents} cảnh báo ba thách thức lớn:
\begin{itemize}
  \item \textbf{Vòng lặp vô hạn (Infinite Loops):} tác tử gọi công cụ lặp đi lặp lại khi không lập được kế hoạch hoàn
        chỉnh hoặc không đánh giá được kết quả.
  \item \textbf{Thất bại lan truyền (Cascading Failures):} trong hệ đa tác tử dựng trên cùng một mô hình nền tảng, một
        điểm yếu chung hoặc dữ liệu sai từ một tác tử có thể kéo theo cả chuỗi.
  \item \textbf{Rò rỉ đặc quyền (Privilege Escalation):} tác tử được cấp quyền truy cập hệ thống có thể bị lợi dụng, ví dụ
        bằng tấn công chèn lệnh (prompt injection), để thực hiện thao tác vượt thẩm quyền.
\end{itemize}

Để bảo đảm an toàn, IBM đề xuất bốn nguyên tắc quản trị:
\begin{enumerate}
  \item \textbf{Audit Trail} -- nhật ký hoạt động có dấu thời gian cho mọi lời gọi công cụ và quyết định, để truy vết
        khi có sự cố.
  \item \textbf{HITL (Human-in-the-loop)} -- con người có thể can thiệp, phê duyệt hoặc dừng khẩn cấp, đặc biệt với
        thao tác có hệ quả lớn.
  \item \textbf{Unique Agent IDs} -- mỗi tác tử có định danh duy nhất để biết hành động nào do tác tử nào thực hiện.
  \item \textbf{Guardrails} -- ngưỡng giới hạn cứng như số vòng lặp tối đa (\texttt{max\_iterations}), thời gian chờ,
        cơ chế tự ngắt.
\end{enumerate}

\hinhtikz{ch6_governance}{Ba nhóm rủi ro vận hành và bốn nguyên tắc quản trị tương ứng.}{ch6-governance}

\begin{table}[!htbp]\centering
\caption{Ma trận kiểm soát rủi ro vận hành AI Agent theo khuyến nghị của IBM và Google Cloud.}\label{tab:ch6-risk}
\small
\begin{tabularx}{\linewidth}{L{3.1cm}XX}
\toprule
\textbf{Nhóm rủi ro} & \textbf{Biểu hiện và hậu quả} & \textbf{Cơ chế kiểm soát} \\
\midrule
Vòng lặp vô hạn & Tác tử gọi công cụ lặp lại khi gặp lỗi, cạn kiệt tài nguyên máy chủ và chi phí token &
Ngưỡng cứng \texttt{max\_iterations}, giới hạn thời gian, tự ngắt \\
Thất bại dây chuyền & Dữ liệu sai từ một tác tử lan sang chuỗi tác tử kế tiếp trong mạng mesh & Định danh tác tử,
kiểm tra dữ liệu trung gian, cô lập lỗi cục bộ \\
Rò rỉ quyền hạn & Tác tử bị chèn lệnh, thực hiện thao tác phá hoại dữ liệu hoặc tài khoản thanh toán & HITL cho thao
tác quan trọng, quyền tối thiểu (least privilege) \\
\bottomrule
\end{tabularx}
\end{table}

Tác tử Q-Learning của bài thực hành cũng có một guardrail theo đúng tinh thần này: hằng số \texttt{MAX\_STEPS} giới hạn
số bước tối đa của mỗi episode, nên ở giai đoạn đầu, khi tác tử còn đi ngẫu nhiên, nó không thể lặp vô hạn trong mê cung.
